\section{Limitations}
\label{sec:limitations}

Loss Susceptibility is local to $\theta_0$, restricted to the declared
parameter block, and dependent on the parameterization and loss. Its
estimate additionally depends on precision, perturbation radius, and the finite
direction bank. It does not encode update sign, realized alignment, curvature,
or off-block damage. Estimating it with $\widehat q_G$ uses $2R$ pool-forward
sweeps that avoid candidate-side
backpropagation during target profiling, but do not guarantee fewer FLOPs or
lower wall-clock time. The exact-rank certificate is evaluated on one declared
setting-level fidelity support; target candidate supports receive separate
split-bank, perturbation-survival, and integrity checks rather than their own
$q_G^\star$ comparison. Representation Proximity measures answer-token
hidden-state similarity, not causal dependence.

The combined score is not fully calibration-free. Its scalar weight is fitted
from completed, target-request-disjoint development trajectories and should be
revalidated when the model, request type, trainable block, or parent unlearner
changes. The method also assumes an enumerated prompt--answer universe. Its
primary outcome, $\Delta$NLL, is fine-grained and scalable but cannot capture
every change in generation, refusal, calibration, or downstream task
performance.

Repaired selectors share a feasibility region rather than identical final
forgetting scores, so forgetting and utility slack must accompany every
comparison. Frozen guards evaluate extraction, generation, and utility only
on declared sets; they do not establish broad resistance to unenumerated
prompts, relearning, or adversarial extraction. The framework is a
prospective screening and allocation layer, neither a new unlearning objective
nor a deletion certificate.
